% !TeX root = ../papers/nc-p3.tex
\section{可変跡の円板での中心四次式と一次経路の積}
\label{sec:rank-fifteen-first-center-products}
\begin{proposition}
定理\ref{thm:rank-fifteen-variable-zero-disk}の領域
$0<\nu(y)<3\nu(x)$ で、同定理で選んだ四中心式を $c_1,\ldots,c_4$ とする。
各始点において、各 $c_j$ に終端から出る四つの一次経路を掛けて得る
十六本は一次独立である。従ってこの四中心商の長さ5の文字別次元は
$(10,10,10,10)$、総次元は40である。
\end{proposition}
\begin{proof}
前節と同じ局所環 $R$ と変数 $H$ を用いる。精度9から長さ5までの未切断経路を
同じ順序で消去すると、長さ4と5で二桁ずつを使い、精度5が残る。
中心方程式の主元は単元なので、前節と同じ自由座標で四中心式を取れる。
各文字 $h$ に対し、文字 $g$ の中心式と末尾の文字 $g\oplus h$ の積を一行とする。
各単語を未切断経路の長さ5の基底で表すと、始点と $h$ の各組に対し
四行十四列の行列が得られる。

\path{computations/verify_rank_fifteen_first_center_products.py}
はこの十六行列を $R$ 上の同じ単元判定で消去する。
各行列とも、階数4に達するまでに必要な $H$ 除算は四回であり、精度が一桁残る。
各除算で捨てる定数項は $\mathbb Q(z)$ で恒等的に零であることを検査する。
従って $H$ が正付値で非零の全ての窓で、十六本の積が文字ごとに四本ずつ独立となる。
中心性により一次経路を先頭に置いた積も同じ空間に属し、
四中心式が生成する両側イデアルの長さ5成分をこれで尽くす。
\end{proof}
